\subsection{A finite definable set generating an infinite subgroup: Problem 21.106}
\label{subsec:concise}
\problemstatement{21.106}

Following Conte and Petschick~\cite{contepetschick}, a formula
$\varphi(x)$ is \emph{concise} in a class of groups if every finite
truth set in that class generates a finite subgroup.
Petschick's \kn{21.106} asks whether every first-order formula with
one free variable is concise in residually finite groups.
Here formulas may have alternating quantifiers; restricting to
group words or to existential formulas gives different questions.

\begin{theorem}
\label{thm:concise}
There is a parameter-free first-order formula whose truth set in
the integral Heisenberg group is $\{z,z^{-1}\}$, where $z$ generates
its infinite cyclic center. In particular, the formula is not
concise in residually finite groups.
\end{theorem}

\begin{proof}
Use the following abbreviations in the language of groups:
\begin{align*}
\operatorname{Cent}(z)&:\quad \forall t\;[z,t]=1,\\
P(u)&:\quad \forall z\,
 \bigl(\operatorname{Cent}(z)\Longrightarrow
                  \exists t\;[u,t]=z\bigr),\\
Q(u,v)&:\quad \forall g\,\exists h\,\exists k\,
 \bigl(g=hk\ \wedge\ [h,u]=1\ \wedge\ [k,v]=1\bigr).
\end{align*}
Define
\[
\varphi(x):\quad \exists u\,\exists v\,
 \bigl(x=[u,v]\ \wedge\ P(u)\ \wedge\ P(v)\ \wedge\ Q(u,v)\bigr).
\]
After renaming bound variables and expanding the abbreviations,
this is a parameter-free formula with exactly one free variable.
All its quantifiers range over group elements.

Write the integral Heisenberg group as
\[
H=\Z^3,\qquad
(a,b,c)(d,e,f)=(a+d,b+e,c+f+ae).
\]
Matrix multiplication, or direct calculation, gives
\[
(a,b,c)^{-1}=(-a,-b,ab-c),\qquad
[(a,b,c),(d,e,f)]=(0,0,ae-db).
\]
Thus $Z(H)=\{(0,0,c):c\in\Z\}=\langle z\rangle$ with
$z=(0,0,1)$. Reduction of coordinates modulo $m$ is a homomorphism
to a finite group of order $m^3$. Choosing $m$ larger than the
absolute value of a nonzero coordinate separates any nonidentity
element, so $H$ is residually finite.

Let $\pi:H\to\Z^2$ be projection onto the first two coordinates.
For $u=(a,b,c)$, the third coordinates of $[u,t]$ form the ideal
$a\Z+b\Z$. Consequently
\[
H\models P(u)\quad\Longleftrightarrow\quad \gcd(a,b)=1.
\]
If $(a,b)$ is primitive, the solutions of $as-br=0$ in
$(r,s)\in\Z^2$ are precisely the multiples of $(a,b)$.
For example, choose $\alpha a+\beta b=1$ and put
$k=\alpha r+\beta s$; the determinant relation gives
$ka=r$ and $kb=s$. Hence
\[
\pi(C_H(u))=\Z(a,b),
\]
and the centralizer is the full inverse image of this subgroup.

Suppose $P(u)$ and $P(v)$ hold, with
$\pi(u)=(a,b)$ and $\pi(v)=(d,e)$.
Both centralizers contain $Z(H)$, so their set product is all of
$H$ exactly when their projected subgroups sum to $\Z^2$.
Indeed, a factorization after projection can be lifted, and its
central error absorbed into the second factor. It follows that
\[
H\models Q(u,v)
\quad\Longleftrightarrow\quad
\Z(a,b)+\Z(d,e)=\Z^2
\quad\Longleftrightarrow\quad ae-db=\pm1.
\]
Thus any value satisfying $\varphi$ is $z$ or $z^{-1}$.
Conversely, the ordered pairs
$((1,0,0),(0,1,0))$ and $((0,1,0),(1,0,0))$
witness the two values. Therefore
$H_\varphi=\{z,z^{-1}\}$, but
$\langle H_\varphi\rangle\cong\Z$.
\end{proof}

The counterexample is already two-generated, torsion-free, and
nilpotent of class two. It is residually a finite $p$-group for
each prime $p$, by reduction modulo powers of $p$.
The elements used to describe witnesses do not occur as parameters
in the formula. No assertion that first-order truth commutes with
finite quotients is needed. The existential-formula theorem of
Conte and Petschick does not apply to this alternating formula.
The experimental controls checked commutator arithmetic and finite
interpretations, while the proof above establishes the infinite
statement.

Contemporaneous work recorded in a public repository credits Lily Zhang
with a counterexample and proof and Evan Li with a Heisenberg
simplification; it also supplies a Lean formalisation\cite{zhang-li-concise}.
\iffullpaper
Our first local proof commit is dated 10 September 2026, 21:05:23 UTC;
the commit adding those Lean files is dated 11 September, 10:54:21 UTC.
These are artifact timestamps, not discovery or verified upload times,
and do not establish priority. We have not built that Lean development.
The retrieval index in \reproductionlocation{} provides a limited
record of source access, not a test of independence.
\else
These related constructions are credited without a claim of discovery
priority. The reported Lean formalisation has not been built here.
\fi

\paragraph{Version 2 attribution.}
Jayadevan~\cite{jayadevan-21-106} gives a different positive formula
with the same two-element truth set in the integral Heisenberg group.
Over any commutative ring it defines the central units. This is
additional to the Zhang--Li work acknowledged above; see
Appendix~\ref{app:contemporary-comparison} for the comparison and chronology.


